% Options for packages loaded elsewhere
\PassOptionsToPackage{unicode}{hyperref}
\PassOptionsToPackage{hyphens}{url}
%
\documentclass[
  12pt,
  a4paper,
  DIV=11,
  numbers=noendperiod]{scrreprt}

\usepackage{amsmath,amssymb}
\usepackage{setspace}
\usepackage{iftex}
\ifPDFTeX
  \usepackage[T1]{fontenc}
  \usepackage[utf8]{inputenc}
  \usepackage{textcomp} % provide euro and other symbols
\else % if luatex or xetex
  \usepackage{unicode-math}
  \defaultfontfeatures{Scale=MatchLowercase}
  \defaultfontfeatures[\rmfamily]{Ligatures=TeX,Scale=1}
\fi
\usepackage{lmodern}
\ifPDFTeX\else  
    % xetex/luatex font selection
\fi
% Use upquote if available, for straight quotes in verbatim environments
\IfFileExists{upquote.sty}{\usepackage{upquote}}{}
\IfFileExists{microtype.sty}{% use microtype if available
  \usepackage[]{microtype}
  \UseMicrotypeSet[protrusion]{basicmath} % disable protrusion for tt fonts
}{}
\usepackage{xcolor}
\setlength{\emergencystretch}{3em} % prevent overfull lines
\setcounter{secnumdepth}{5}
% Make \paragraph and \subparagraph free-standing
\ifx\paragraph\undefined\else
  \let\oldparagraph\paragraph
  \renewcommand{\paragraph}[1]{\oldparagraph{#1}\mbox{}}
\fi
\ifx\subparagraph\undefined\else
  \let\oldsubparagraph\subparagraph
  \renewcommand{\subparagraph}[1]{\oldsubparagraph{#1}\mbox{}}
\fi


\providecommand{\tightlist}{%
  \setlength{\itemsep}{0pt}\setlength{\parskip}{0pt}}\usepackage{longtable,booktabs,array}
\usepackage{calc} % for calculating minipage widths
% Correct order of tables after \paragraph or \subparagraph
\usepackage{etoolbox}
\makeatletter
\patchcmd\longtable{\par}{\if@noskipsec\mbox{}\fi\par}{}{}
\makeatother
% Allow footnotes in longtable head/foot
\IfFileExists{footnotehyper.sty}{\usepackage{footnotehyper}}{\usepackage{footnote}}
\makesavenoteenv{longtable}
\usepackage{graphicx}
\makeatletter
\def\maxwidth{\ifdim\Gin@nat@width>\linewidth\linewidth\else\Gin@nat@width\fi}
\def\maxheight{\ifdim\Gin@nat@height>\textheight\textheight\else\Gin@nat@height\fi}
\makeatother
% Scale images if necessary, so that they will not overflow the page
% margins by default, and it is still possible to overwrite the defaults
% using explicit options in \includegraphics[width, height, ...]{}
\setkeys{Gin}{width=\maxwidth,height=\maxheight,keepaspectratio}
% Set default figure placement to htbp
\makeatletter
\def\fps@figure{htbp}
\makeatother

\usepackage{indentfirst}
\usepackage{float}
\floatplacement{figure}{H}
\IfFileExists{plex-otf.sty}{
  %% Full TeXlive
  \usepackage[
    % math,
    RM={Scale=0.94},SS={Scale=0.94},SScon={Scale=0.94},TT={Scale=MatchLowercase,FakeStretch=0.9},
    DefaultFeatures={Ligatures=Common}
  ]{plex-otf}
  % \usepackage[Scale=MatchUppercase]{juliamono}
}{
  %% TinyTeX
  \usepackage{libertine}
}
\KOMAoption{captions}{tableheading}
\makeatletter
\@ifpackageloaded{caption}{}{\usepackage{caption}}
\AtBeginDocument{%
\ifdefined\contentsname
  \renewcommand*\contentsname{Содержание}
\else
  \newcommand\contentsname{Содержание}
\fi
\ifdefined\listfigurename
  \renewcommand*\listfigurename{Список иллюстраций}
\else
  \newcommand\listfigurename{Список иллюстраций}
\fi
\ifdefined\listtablename
  \renewcommand*\listtablename{Список таблиц}
\else
  \newcommand\listtablename{Список таблиц}
\fi
\ifdefined\figurename
  \renewcommand*\figurename{Рисунок}
\else
  \newcommand\figurename{Рисунок}
\fi
\ifdefined\tablename
  \renewcommand*\tablename{Таблица}
\else
  \newcommand\tablename{Таблица}
\fi
}
\@ifpackageloaded{float}{}{\usepackage{float}}
\floatstyle{ruled}
\@ifundefined{c@chapter}{\newfloat{codelisting}{h}{lop}}{\newfloat{codelisting}{h}{lop}[chapter]}
\floatname{codelisting}{Список}
\newcommand*\listoflistings{\listof{codelisting}{Листинги}}
\makeatother
\makeatletter
\makeatother
\makeatletter
\@ifpackageloaded{caption}{}{\usepackage{caption}}
\@ifpackageloaded{subcaption}{}{\usepackage{subcaption}}
\makeatother
\ifLuaTeX
\usepackage[bidi=basic]{babel}
\else
\usepackage[bidi=default]{babel}
\fi
\babelprovide[main,import]{russian}
\babelprovide[import]{english}
% get rid of language-specific shorthands (see #6817):
\let\LanguageShortHands\languageshorthands
\def\languageshorthands#1{}
\ifLuaTeX
  \usepackage{selnolig}  % disable illegal ligatures
\fi
\usepackage[style=gost-numeric,backend=biber,langhook=extras,autolang=other*]{biblatex}
\addbibresource{bib/cite.bib}
\usepackage{csquotes}
\usepackage{bookmark}

\IfFileExists{xurl.sty}{\usepackage{xurl}}{} % add URL line breaks if available
\urlstyle{same} % disable monospaced font for URLs
\hypersetup{
  pdftitle={Групповой проект. Этап 2},
  pdfauthor={Вакутайпа Милдред, Разанацуа Сара Естэлл, Нджову Нелиа},
  pdflang={ru-RU},
  hidelinks,
  pdfcreator={LaTeX via pandoc}}

\title{Групповой проект. Этап 2}
\usepackage{etoolbox}
\makeatletter
\providecommand{\subtitle}[1]{% add subtitle to \maketitle
  \apptocmd{\@title}{\par {\large #1 \par}}{}{}
}
\makeatother
\subtitle{Алгоритм решения задачи Рост дендритов}
\author{Вакутайпа Милдред, Разанацуа Сара Естэлл, Нджову Нелиа}
\date{}

\begin{document}
\maketitle

\renewcommand*\contentsname{Содержание}
{
\setcounter{tocdepth}{1}
\tableofcontents
}
\listoffigures
\listoftables
\setstretch{1.5}
\chapter{Введение}\label{ux432ux432ux435ux434ux435ux43dux438ux435}

\section{Актуальность}\label{ux430ux43aux442ux443ux430ux43bux44cux43dux43eux441ux442ux44c}

Дендрит - это характерная древовидная структура кристаллов, растущих по
мере затвердевания расплавленного металла, форма, получаемая в
результате более быстрого роста в энергетически выгодных
кристаллографических направлениях.

Процессы кристаллизации и роста дендритов играют ключевую роль в
различных областях науки и техники. В металлургии и материаловедении
структура затвердевших металлов и сплавов, образующаяся в процессе
литья, сварки и аддитивного производства, определяет их механические
свойства --- прочность, пластичность и усталостную долговечность.
Дендритная микроструктура, возникающая при неравновесном затвердевании,
непосредственно влияет на распределение примесей, образование пор и
трещин, а также на последующую термообработку материала.

В природе дендритные структуры также широко распространены: от снежинок
и морозных узоров на стекле до минеральных образований в геологических
процессах. Понимание механизмов роста дендритов важно не только для
технологических применений, но и для фундаментальной физики
неравновесных процессов структурообразования. Несмотря на многолетние
исследования, проблема управляемого управления дендритной
микроструктурой остается актуальной.

\section{Цель}\label{ux446ux435ux43bux44c}

Целью данной работы является написание алгоритма реализации модели роста
дендритов.

\section{Задачи}\label{ux437ux430ux434ux430ux447ux438}

Изучите алгоритм модели роста дендритов и опишите основные этапы работы
алгоритма

\chapter{Алгоритм}\label{ux430ux43bux433ux43eux440ux438ux442ux43c}

\section{Шаг 1: Настройка параметров
модели}\label{ux448ux430ux433-1-ux43dux430ux441ux442ux440ux43eux439ux43aux430-ux43fux430ux440ux430ux43cux435ux442ux440ux43eux432-ux43cux43eux434ux435ux43bux438}

Первый этап включает в себя определение физических свойств материала и
вычислительных параметров для моделирования. Корректность входных данных
определяет надежность всего процесса моделирования.

\subsection{Физические свойства
материала:}\label{ux444ux438ux437ux438ux447ux435ux441ux43aux438ux435-ux441ux432ux43eux439ux441ux442ux432ux430-ux43cux430ux442ux435ux440ux438ux430ux43bux430}

\begin{itemize}
\tightlist
\item
  Плотность \(\rho\): определяет массу на единицу объема материала.
\item
  Удельная теплоемкость \(c_p\): количество энергии, необходимо для
  нагрева единицы массы на один градус.
\item
  Скрытая теплота плавления \(L\): энергия фазавого перехода
\item
  Теплопроводность \(\kappa\): способность материала передавать тепловую
  энергию
\item
  Коэффициент температуропроводности \(\chi = \kappa/(\rho c_p)\):
  скорость выравнивания температурного градиента
\item
  Температура плавления \(T_m\): критическая точка фазового перехода
\item
  Коэффициент поверхностного натяжения \(\gamma\): влияет на
  стабильность границы раздела фаз
\item
  Кинетический коэффициент \(\beta\): характеризует скорость
  прикрепления атомов к поверхности кристалла{[}1,2{]}
\end{itemize}

\subsection{Вычислительные
параметры}\label{ux432ux44bux447ux438ux441ux43bux438ux442ux435ux43bux44cux43dux44bux435-ux43fux430ux440ux430ux43cux435ux442ux440ux44b}

\begin{itemize}
\tightlist
\item
  Размер сетки \(N \times N\): Количество узлов по каждому измерению
\item
  Пространственный шаг \(h\): Расстояние между соседними узлами
\item
  Временной шаг роста \(\Delta\): Интервал одного шага роста
\item
  Число подшагов \(m\): Количество подшагов тепловой диффузии
\item
  Капиллярный параметр \(\lambda\): Интенсивность стабилизации границы
\item
  Амплитуда шума \(\delta\): Уровень тепловых флуктуаций{[}3{]}
\end{itemize}

\subsection{Начальные и граничные
условия:}\label{ux43dux430ux447ux430ux43bux44cux43dux44bux435-ux438-ux433ux440ux430ux43dux438ux447ux43dux44bux435-ux443ux441ux43bux43eux432ux438ux44f}

\begin{itemize}
\tightlist
\item
  Начальная температура расплава \(T_\infty\): определяет степень
  переохлаждения системы
\item
  Безразмерное переохлаждение \(S = c_p (T_m - T_\infty)/L\): ключевой
  параметр термодинамической готовности системы
\item
  Граничные условия: постоянная температура на всех границах расчетной
  области, равная начальной температуре расплава \(T_\infty\){[}4{]}
\end{itemize}

\section{Шаг 2: Инициализация структур
данных}\label{ux448ux430ux433-2-ux438ux43dux438ux446ux438ux430ux43bux438ux437ux430ux446ux438ux44f-ux441ux442ux440ux443ux43aux442ux443ux440-ux434ux430ux43dux43dux44bux445}

На втором этапе создаются структуры данных для хранения состояния
системы и задается начальная конфигурация.

\subsection{Структуры
данных:}\label{ux441ux442ux440ux443ux43aux442ux443ux440ux44b-ux434ux430ux43dux43dux44bux445}

Создаются два двумерных массива размером \(N \times N\):

\begin{enumerate}
\def\labelenumi{\arabic{enumi}.}
\tightlist
\item
  Массив фазового состояния \(phase[i][j]\):
\end{enumerate}

\begin{itemize}
\tightlist
\item
  \(phase[i][j] = 1\) --- твердая фаза
\item
  \(phase[i][j] = 0\) --- жидкая фаза
\end{itemize}

\begin{enumerate}
\def\labelenumi{\arabic{enumi}.}
\setcounter{enumi}{1}
\tightlist
\item
  Массив безразмерной температуры \(T[i][j]\):
\end{enumerate}

\begin{itemize}
\tightlist
\item
  Хранит текущее значение безразмерной температуры в каждом узле
\end{itemize}

\subsection{Инициализация:}\label{ux438ux43dux438ux446ux438ux430ux43bux438ux437ux430ux446ux438ux44f}

\begin{enumerate}
\def\labelenumi{\arabic{enumi}.}
\tightlist
\item
  Заполнение массивов начальными значениями:
\end{enumerate}

\begin{itemize}
\tightlist
\item
  Все узлы: \(phase[i][j] = 0\) (жидкое состояние)
\item
  Все узлы: \(T[i][j] = 0\) (безразмерная температура расплава)
\end{itemize}

\begin{enumerate}
\def\labelenumi{\arabic{enumi}.}
\setcounter{enumi}{1}
\tightlist
\item
  Создание затравки кристаллизации:
\end{enumerate}

\begin{itemize}
\tightlist
\item
  В центральной области сетки задается твердая затравка (например,
  квадрат \(5 \times 5\) узлов)
\item
  Для узлов затравки: \(phase[i][j] = 1\) (твердое состояние)
\item
  Для узлов затравки: \(T[i][j] = \tilde{T}_m\) (безразмерная
  температура плавления)
\end{itemize}

\section{Шаг 3: Расчет температурного
поля}\label{ux448ux430ux433-3-ux440ux430ux441ux447ux435ux442-ux442ux435ux43cux43fux435ux440ux430ux442ux443ux440ux43dux43eux433ux43e-ux43fux43eux43bux44f}

Третий шаг реализует численное решение уравнения теплопроводности для
расчета эволюции температурного поля.{[}3{]}

\subsection{Дискретизация уравнения
теплопроводности:}\label{ux434ux438ux441ux43aux440ux435ux442ux438ux437ux430ux446ux438ux44f-ux443ux440ux430ux432ux43dux435ux43dux438ux44f-ux442ux435ux43fux43bux43eux43fux440ux43eux432ux43eux434ux43dux43eux441ux442ux438}

Используется явная конечно-разностная схема. Лапласиан в узле \((i,j)\)
аппроксимируется с учетом ближайших и диагональных соседей:

\[
\nabla^2 T \approx \frac{\langle T_{(i,j)} \rangle - T_{i,j}}{(4 + 4w)(1 + 2w)h^2}
\]

где среднее значение температуры соседей:

\[
\langle T_{(i,j)} \rangle = \frac{ (T_{i+1,j} + T_{i-1,j} + T_{i,j+1} + T_{i,j-1}) + w (T_{i+1,j+1} + T_{i+1,j-1} + T_{i-1,j+1} + T_{i-1,j-1}) }{4 + 4w}
\]

Коэффициент \(w = 1/2\) определяет вклад диагональных соседей.{[}4{]}

\subsection{Алгоритм расчета тепловой
диффузии:}\label{ux430ux43bux433ux43eux440ux438ux442ux43c-ux440ux430ux441ux447ux435ux442ux430-ux442ux435ux43fux43bux43eux432ux43eux439-ux434ux438ux444ux444ux443ux437ux438ux438}

Для каждого шага роста \(\Delta t\) выполняется \(m\) подшагов тепловой
диффузии:

\begin{enumerate}
\def\labelenumi{\arabic{enumi}.}
\item
  Для каждого узла сетки \((i, j)\) вычислить лапласиан
  \(\nabla^2 T_{i,j}\) по приведенной формуле
\item
  Обновить температуру:
  \(T_{i,j}^{new} = T_{i,j} + \chi \frac{\Delta t}{m} \nabla^2 T_{i,j}\)
\item
  Применить граничные условия постоянной температуры:
  \(T_{граница} = 0\)
\end{enumerate}

\subsection{Условие
устойчивости:}\label{ux443ux441ux43bux43eux432ux438ux435-ux443ux441ux442ux43eux439ux447ux438ux432ux43eux441ux442ux438}

\[
\frac{\chi \Delta t}{m h^2} < \frac{1}{4}
\]

\section{Шаг 4: Моделирование роста
дендритов}\label{ux448ux430ux433-4-ux43cux43eux434ux435ux43bux438ux440ux43eux432ux430ux43dux438ux435-ux440ux43eux441ux442ux430-ux434ux435ux43dux434ux440ux438ux442ux43eux432}

На четвертом этапе после завершения всех подшагов тепловой диффузии
выполняется проверка условий затвердевания и обновление фазового
состояния.

\subsection{Порядок
обновления:}\label{ux43fux43eux440ux44fux434ux43eux43a-ux43eux431ux43dux43eux432ux43bux435ux43dux438ux44f}

Важной особенностью алгоритма является последовательное обновление:

\begin{enumerate}
\def\labelenumi{\arabic{enumi}.}
\item
  Сначала полностью рассчитывается температурное поле на текущем шаге
  (все \(m\) подшагов)
\item
  Затем, на основе полученного температурного поля, проверяются условия
  затвердевания
\end{enumerate}

\subsection{Вычисление кривизны
границы:}\label{ux432ux44bux447ux438ux441ux43bux435ux43dux438ux435-ux43aux440ux438ux432ux438ux437ux43dux44b-ux433ux440ux430ux43dux438ux446ux44b}

Для жидкого узла, имеющего хотя бы один твердый сосед, вычисляется
дискретизированная кривизна:{[}3{]}

\[
s_{i,j} = \sum_{\text{nearest}} n_{i,j} + w_n \sum_{\text{diagonal}} n_{i,j} - \left( \frac{5}{2} + \frac{5}{2} w_n \right)
\]

где:

\begin{itemize}
\tightlist
\item
  \(n_{i,j} = 1\) если соседний узел твердый, и \(0\) в противном случае
\item
  \(w_n = 1/2\) --- весовой коэффициент для диагональных соседей
\item
  \(s_{i,j} < 0\) для выпуклых участков (вершин), \(s_{i,j} > 0\) для
  вогнутых
\end{itemize}

\subsection{Критерий
затвердевания:}\label{ux43aux440ux438ux442ux435ux440ux438ux439-ux437ux430ux442ux432ux435ux440ux434ux435ux432ux430ux43dux438ux44f}

Жидкий узел переходит в твердое состояние, если выполняется условие:

\[
T_{i,j} \le \tilde{T}_m + \lambda s_{i,j} + \eta_{i,j} \delta
\]

где:

\begin{itemize}
\item
  \(\tilde{T}_m\) --- безразмерная температура плавления
\item
  \(\lambda\) --- капиллярный параметр, объединяющий эффекты
  Гиббса--Томсона и кинетического переохлаждения{[}2{]}
\item
  \(s_{i,j}\) --- дискретизированная кривизна
\item
  \(\eta_{i,j}\) --- случайное число из равномерного распределения
  \([-1, 1]\)
\item
  \(\delta\) --- амплитуда теплового шума{[}3{]}
\end{itemize}

\subsection{Алгоритм проверки
затвердевания:}\label{ux430ux43bux433ux43eux440ux438ux442ux43c-ux43fux440ux43eux432ux435ux440ux43aux438-ux437ux430ux442ux432ux435ux440ux434ux435ux432ux430ux43dux438ux44f}

\begin{enumerate}
\def\labelenumi{\arabic{enumi}.}
\item
  Полный обход всех узлов сетки \((i, j)\)
\item
  Для каждого узла:
\end{enumerate}

\begin{itemize}
\tightlist
\item
  Если \(phase[i][j] = 1\) (твердый), пропустить
\item
  Если \(phase[i][j] = 0\) (жидкий), проверить наличие твердых соседей
\item
  Если твердые соседи отсутствуют, пропустить
\item
  Если твердые соседи есть:

  \begin{itemize}
  \tightlist
  \item
    Вычислить \(s_{i,j}\) на основе соседних узлов
  \item
    Сгенерировать случайное число \(\eta_{i,j}\)
  \item
    Проверить критерий затвердевания
  \item
    Если условие выполнено:

    \begin{itemize}
    \tightlist
    \item
      \(phase[i][j] = 1\) (переход в твердое состояние)
    \item
      \(T[i][j] = T[i][j] + 1\) (выделение скрытой теплоты){[}3{]}
    \end{itemize}
  \end{itemize}
\end{itemize}

\section{Шаг 5: Организация основного цикла
вычислений}\label{ux448ux430ux433-5-ux43eux440ux433ux430ux43dux438ux437ux430ux446ux438ux44f-ux43eux441ux43dux43eux432ux43dux43eux433ux43e-ux446ux438ux43aux43bux430-ux432ux44bux447ux438ux441ux43bux435ux43dux438ux439}

Пятый этап описывает структуру основного вычислительного цикла.

\subsection{Блок-схема
алгоритма:}\label{ux431ux43bux43eux43a-ux441ux445ux435ux43cux430-ux430ux43bux433ux43eux440ux438ux442ux43cux430}

Алгоритм моделирования роста дендритов состоит из следующих основных
этапов:

\begin{enumerate}
\def\labelenumi{\arabic{enumi}.}
\item
  Начало: Запуск программы
\item
  Шаг 1: Настройка параметров модели
\end{enumerate}

\begin{itemize}
\tightlist
\item
  Определение физических свойств материала (ρ, c\_p, κ, L, T\_m, γ, β)
\item
  Определение вычислительных параметров (N, h, Δt, m, λ, δ)
\item
  Определение начальных и граничных условий (T\_∞, T\_граница = T\_∞)
\end{itemize}

\begin{enumerate}
\def\labelenumi{\arabic{enumi}.}
\setcounter{enumi}{2}
\tightlist
\item
  Шаг 2: Инициализация массивов
\end{enumerate}

\begin{itemize}
\tightlist
\item
  Создание массива фазового состояния phase{[}N{]}{[}N{]} (все узлы = 0
  --- жидкие)
\item
  Создание массива температуры T{[}N{]}{[}N{]} (все узлы = 0)
\item
  Создание твердой затравки в центре (phase = 1, T = T̃\_m)
\end{itemize}

\begin{enumerate}
\def\labelenumi{\arabic{enumi}.}
\setcounter{enumi}{3}
\tightlist
\item
  Проверка условия остановки
\end{enumerate}

\begin{itemize}
\tightlist
\item
  Если дендрит достиг границы → переход к завершению
\item
  Если не достиг → продолжение вычислений
\end{itemize}

\begin{enumerate}
\def\labelenumi{\arabic{enumi}.}
\setcounter{enumi}{4}
\tightlist
\item
  Шаг 3: Расчет температурного поля
\end{enumerate}

\begin{itemize}
\tightlist
\item
  Выполнение m подшагов тепловой диффузии:

  \begin{itemize}
  \tightlist
  \item
    Для каждого узла вычислить лапласиан ∇²T
  \item
    Обновить температуру T\_new = T + χ·(Δt/m)·∇²T
  \item
    Применить граничные условия T\_граница = 0
  \end{itemize}
\end{itemize}

\begin{enumerate}
\def\labelenumi{\arabic{enumi}.}
\setcounter{enumi}{5}
\tightlist
\item
  Шаг 4: Моделирование роста
\end{enumerate}

\begin{itemize}
\tightlist
\item
  Обход всех узлов сетки:

  \begin{itemize}
  \tightlist
  \item
    Для каждого жидкого узла с твердыми соседями:

    \begin{itemize}
    \tightlist
    \item
      Вычислить кривизну s\_ij
    \item
      Сгенерировать случайный шум η\_ij
    \item
      Проверить критерий затвердевания: T ≤ T̃\_m + λ·s\_ij + η·δ
    \item
      Если условие выполнено:

      \begin{itemize}
      \tightlist
      \item
        phase{[}i{]}{[}j{]} = 1 (переход в твердое состояние)
      \item
        T{[}i{]}{[}j{]} = T{[}i{]}{[}j{]} + 1 (выделение скрытой
        теплоты)
      \end{itemize}
    \end{itemize}
  \end{itemize}
\end{itemize}

\begin{enumerate}
\def\labelenumi{\arabic{enumi}.}
\setcounter{enumi}{6}
\tightlist
\item
  Сохранение результатов
\end{enumerate}

\begin{itemize}
\tightlist
\item
  Запись фазового и температурного полей для визуализации
\end{itemize}

\begin{enumerate}
\def\labelenumi{\arabic{enumi}.}
\setcounter{enumi}{7}
\item
  Возврат к проверке условия остановки (пункт 4)
\item
  Конец: Завершение моделирования
\end{enumerate}

\subsection{Условие
остановки:}\label{ux443ux441ux43bux43eux432ux438ux435-ux43eux441ux442ux430ux43dux43eux432ux43aux438}

Вычисления продолжаются до тех пор, пока дендрит не достигнет границы
расчетной области. Условие проверяется после каждого шага роста: если
любой твердый узел находится на границе сетки, моделирование
завершается.

\section{Шаг 6: Исследование зависимости характеристик роста от
параметров}\label{ux448ux430ux433-6-ux438ux441ux441ux43bux435ux434ux43eux432ux430ux43dux438ux435-ux437ux430ux432ux438ux441ux438ux43cux43eux441ux442ux438-ux445ux430ux440ux430ux43aux442ux435ux440ux438ux441ux442ux438ux43a-ux440ux43eux441ux442ux430-ux43eux442-ux43fux430ux440ux430ux43cux435ux442ux440ux43eux432}

На шестом этапе выполняются серии расчетов для анализа влияния ключевых
параметров на морфологию дендритов.{[}1{]}

\subsection{Параметры для
исследования:}\label{ux43fux430ux440ux430ux43cux435ux442ux440ux44b-ux434ux43bux44f-ux438ux441ux441ux43bux435ux434ux43eux432ux430ux43dux438ux44f}

\begin{longtable}[]{@{}
  >{\raggedright\arraybackslash}p{(\columnwidth - 4\tabcolsep) * \real{0.2564}}
  >{\raggedright\arraybackslash}p{(\columnwidth - 4\tabcolsep) * \real{0.2564}}
  >{\raggedright\arraybackslash}p{(\columnwidth - 4\tabcolsep) * \real{0.4872}}@{}}
\toprule\noalign{}
\begin{minipage}[b]{\linewidth}\raggedright
Параметр
\end{minipage} & \begin{minipage}[b]{\linewidth}\raggedright
Диапазон
\end{minipage} & \begin{minipage}[b]{\linewidth}\raggedright
Исследуемый эффект
\end{minipage} \\
\midrule\noalign{}
\endhead
\bottomrule\noalign{}
\endlastfoot
Переохлаждение \(\tilde{T}_m\) & 0.1--0.9 & Скорость роста, плотность
ветвления \\
Капиллярный параметр \(\lambda\) & 0.01--0.5 & Стабилизация границы,
радиус вершины \\
Амплитуда шума \(\delta\) & 0--0.2 & Вторичное ветвление, симметрия \\
Температуропроводность \(\chi\) & 0.1--0.5 & Скорость отвода тепла \\
\end{longtable}

\subsection{Анализируемые
характеристики:}\label{ux430ux43dux430ux43bux438ux437ux438ux440ux443ux435ux43cux44bux435-ux445ux430ux440ux430ux43aux442ux435ux440ux438ux441ux442ux438ux43aux438}

\begin{enumerate}
\def\labelenumi{\arabic{enumi}.}
\item
  Скорость роста вершины: измеряется как смещение фронта во времени
\item
  Радиус кривизны вершины: определяется по конфигурации твердой фазы
\item
  Плотность вторичных ветвей: количество ответвлений на единицу длины
\item
  Время достижения границы: полное время моделирования до остановки
\end{enumerate}

\subsection{Методика проведения
исследований:}\label{ux43cux435ux442ux43eux434ux438ux43aux430-ux43fux440ux43eux432ux435ux434ux435ux43dux438ux44f-ux438ux441ux441ux43bux435ux434ux43eux432ux430ux43dux438ux439}

\begin{enumerate}
\def\labelenumi{\arabic{enumi}.}
\item
  Зафиксировать все параметры, кроме исследуемого
\item
  Выполнить серию расчетов с различными значениями исследуемого
  параметра
\item
  Для каждого расчета сохранить:
\end{enumerate}

\begin{itemize}
\tightlist
\item
  Финальную морфологию дендрита
\item
  Зависимость скорости роста от времени
\item
  Время достижения границы
\end{itemize}

\begin{enumerate}
\def\labelenumi{\arabic{enumi}.}
\setcounter{enumi}{3}
\tightlist
\item
  Построить графики зависимостей характеристик от параметров
\end{enumerate}

\section{Шаг 7: Визуализация
результатов}\label{ux448ux430ux433-7-ux432ux438ux437ux443ux430ux43bux438ux437ux430ux446ux438ux44f-ux440ux435ux437ux443ux43bux44cux442ux430ux442ux43eux432}

Седьмой этап обеспечивает графическое представление процесса роста и
финальных результатов(рис. 1)

\begin{figure}

\centering{

\includegraphics[width=0.7\textwidth,height=\textheight]{image/pic1.png}

}

\caption{\label{fig-048}Рост дендрита}

\end{figure}%

\subsection{Визуализация в процессе
моделирования:}\label{ux432ux438ux437ux443ux430ux43bux438ux437ux430ux446ux438ux44f-ux432-ux43fux440ux43eux446ux435ux441ux441ux435-ux43cux43eux434ux435ux43bux438ux440ux43eux432ux430ux43dux438ux44f}

На каждом шаге роста (или с заданной периодичностью) сохраняются:

\begin{enumerate}
\def\labelenumi{\arabic{enumi}.}
\item
  Фазовое поле: отображение твердой (например, черный) и жидкой (белый)
  фаз
\item
  Температурное поле: цветовая карта распределения безразмерной
  температуры
\end{enumerate}

\subsection{Финальная
визуализация:}\label{ux444ux438ux43dux430ux43bux44cux43dux430ux44f-ux432ux438ux437ux443ux430ux43bux438ux437ux430ux446ux438ux44f}

По завершении моделирования:

\begin{itemize}
\tightlist
\item
  Построение финальной морфологии дендрита
\item
  Создание анимации процесса роста на основе сохраненных кадров
\item
  Построение графиков зависимостей (скорость роста от времени, влияние
  параметров)
\end{itemize}

\chapter{Вывод}\label{ux432ux44bux432ux43eux434}

На втором этапе группового проекта описан алгоритм из 7 шагов для
моделирования роста дендритов в переохлажденном расплаве. Граничные
условия задаются постоянной температурой, равной начальной температуре
расплава.

\chapter*{Список
литературы}\label{ux441ux43fux438ux441ux43eux43a-ux43bux438ux442ux435ux440ux430ux442ux443ux440ux44b}
\addcontentsline{toc}{chapter}{Список литературы}

\begin{enumerate}
\def\labelenumi{\arabic{enumi}.}
\item
  W. W. Mullins and R. F. Sekerka, \enquote{Stability of a planar
  interface during solidification of a dilute binary alloy,} Journal of
  Applied Physics, vol.~35, no. 2, pp.~444--451, 1964.
\item
  J. W. Gibbs, \enquote{On the Equilibrium of Heterogeneous Substances,}
  Transactions of the Connecticut Academy of Arts and Sciences, vol.~3,
  pp.~108--248, 1875--1878.
\item
  Y. Saito, G. Goldbeck-Wood, and H. Müller-Krumbhaar,
  \enquote{Numerical simulation of dendritic growth,} Physical Review A,
  vol.~38, no. 4, pp.~2148--2157, 1988.
\item
  Д. А. Медведев, \enquote{Моделирование физических процессов и явлений
  на ПК,} учебное пособие
\item
  Wang K. и др. Dendrite growth in the recharging process of zinc--air
  batteries {[}Электронный ресурс{]}. Journal of Materials Chemistry A,
  2025. URL: https://pubs.rsc.org/en/
  content/articlelanding/2015/ta/c5ta06366c/unauth.
\end{enumerate}


\printbibliography


\end{document}
